\documentclass{article}%
\usepackage[T1]{fontenc}%
\usepackage[utf8]{inputenc}%
\usepackage{lmodern}%
\usepackage{textcomp}%
\usepackage{lastpage}%
\usepackage{geometry}%
\geometry{margin=2.2cm}%
%
\usepackage{amsmath}%
\usepackage{amssymb}%
\usepackage{booktabs}%
\usepackage{longtable}%
\usepackage{array}%
\usepackage{hyperref}%
\hypersetup{hidelinks}%
\title{IME 2027 {-} Matemática\\[4pt]\large Comparação com o gabarito Poliedro\thanks{Gerado com gpt-resolve: \url{https://github.com/lgabs/gpt-resolve}}}%
\date{28/10/2026}%
%
\begin{document}%
\normalsize%
\maketitle%
\section*{Resumo}\textbf{Nota média: 10,0/10} (12 de 12 itens com a resolução correta; cada item vale o mesmo dentro da questão).\par\medskip%
\begin{longtable}{@{}c c p{0.78\textwidth}@{}}\toprule Questão & Nota & Resumo \\ \midrule
1 & 10,0 & A resolução chega a 12$\surd$3 - 18 $\approx$ 2,785, o mesmo valor do gabarito Poliedro, e as duas estão certas. O método é o mesmo (inclusão{-}exclusão com a interseção vista como dodecágono regular de apótema $\surd$3/2), e a resolução ainda confirma o resultado somando as 6 pontas triangulares. Conferi tudo também numericamente, com integração em grade. Não há erro nem ponto duvidoso a levar a um professor. \\
2 & 10,0 & A resolução está correta e coincide com o gabarito Poliedro: P = 2/5. As raízes são cis(k$\pi$/3), k=1..5, e só z = 1/2 $\pm$ ($\surd$3/2)i satisfazem a inequação, com igualdade. Conferi as contas de forma exata e numérica. Não há divergência. \\
3 & 10,0 & A resolução e o gabarito Poliedro chegam à mesma resposta, x = 7, e os dois estão certos. Conferi a resolução passo a passo: as expressões f(x) = 1/2 para x $\geq$ 0 e f(-1/k) = (1/2)(1 - 3\^{}\{1-k\}), a condição 3\^{}\{k-1\} $\geq$ 500 e os valores em x = 6 e x = 7. Também confirmei os resultados com cálculo exato em Python. A resolução é até mais completa que o gabarito, porque elimina os inteiros negativos e x = 0. Não há erros a apontar. \\
4 & 10,0 & A resolução e o gabarito Poliedro dão o mesmo valor, 1345, e os dois estão certos. O raciocínio da resolução está completo: usa a fatoração ?(y\_i+x) = -P(-x), reduz o grau com x$^2$ = 7x + 3 e fecha com as relações de Girard. Conferi as contas de novo, à mão e com cálculo numérico. Não há pontos para levar a um professor. \\
5 & 10,0 & A resolução e o gabarito do Poliedro concordam. Os dois veem que, com o dado 5540 do enunciado (confirmado no PDF oficial), não sai valor racional, e na hipótese de erro de digitação (5440) ambos chegam a q = 2 e S = 10912. A resolução vai além: resolve o caso literal, com S = 512 + 5540t - 512t$^2$ $\approx$ 11196,5 e t $\approx$ 2,5116, e esse valor confere numericamente com as duas condições do enunciado. Não encontrei erro de conta. O único ponto em aberto é a falha do enunciado, que provavelmente leva à anulação da questão. Não consegui gravar a observação sobre o erro do enunciado na memória da equipe: o acesso à pasta de memória foi negado nesta sessão. \\
6 & 10,0 & A resolução está correta e bate com o gabarito Poliedro: P = 1/3. O argumento é completo. A resolução mostra que f\^{}(9)(1) = 1 equivale a o comprimento do ciclo do 1 dividir 9. Também mostra que esse comprimento é uniforme em \{1,\ldots{},9\}, com 8! permutações para cada valor. Testei as 9! permutações por força bruta e encontrei 120960 favoráveis em 362880. Isso confirma 1/3. Não achei divergências nem ambiguidades no enunciado. \\
7 & 10,0 & A resolução e o gabarito batem: os dois chegam a tr(AB) = C(m+n+2, n+1) = C(2026, 1001) e concluem que o traço é divisível por 1013. Conferi numericamente a fórmula e a divisibilidade. A demonstração da resolução é completa e rigorosa: usa Vandermonde, Stifel e a relação 1001$\cdot$C(2026,1001) = 2$\cdot$1013$\cdot$C(2025,1000) com mdc(1001,1013) = 1, e ainda confere por Legendre. O único problema está no gabarito: a frase '1013 não divide 1025!$\cdot$1001!' é falsa, porque 1013 divide 1025!. A conclusão continua certa porque 1013 aparece duas vezes em 2026! e uma só vez no denominador. A resolução merece a nota integral. \\
8 & 10,0 & A resolução está correta e completa, e bate com o gabarito do Poliedro: é o mesmo método (prostaférese, diferença de tangentes, soma telescópica) e chega à mesma identidade. Ela também mostra explicitamente onde cada hipótese é usada, o que o gabarito só resume. Conferi a álgebra e fiz uma verificação numérica em cerca de 20 mil casos aleatórios, sem divergência. Só há dois detalhes de redação, sem impacto na nota. \\
9 & 10,0 & A resolução está correta e coincide com o gabarito Poliedro nos dois itens. O lugar de M é a elipse 4x$^2$/49 + y$^2$ = 1, e os pontos de interseção com x = $\pm$2$\surd$3 são ($\pm$2$\surd$3, $\pm$1/7). O enunciado transcrito confere com o PDF, e refiz e conferi numericamente todas as contas: normal, Q, M, lugar e interseções. A única diferença é que a resolução trata o caso degenerado P=($\pm$4,0), excluindo os vértices ($\pm$7/2, 0), o que o gabarito não faz. Isso não altera a resposta. \\
10 & 10,0 & A resolução e o gabarito Poliedro coincidem nos dois itens, e os dois estão certos. Pelas minhas contas, (a) dá V\_max = 192$\surd$3 cm$^3$ $\approx$ 332,55 cm$^3$, e (b) dá um cubo de aresta 4$\surd$3 cm, ou seja, altura 4$\surd$3 cm e base quadrada de lado 4$\surd$3 cm. A resolução é mais rigorosa que o gabarito: trata uma base quadrilátera qualquer e o prisma apenas contido na bola, e mostra que o máximo é atingido respeitando a hipótese do enunciado. O gabarito supõe uma base retangular inscrita sem justificar e traz um erro de digitação sem importância. \\
\bottomrule\end{longtable}%
\section*{Principais pontos de dúvida}%
\begin{itemize}%
\item%
\textbf{Q5} \textit{(statement)} O enunciado oficial (PDF, pág. 2) traz x2+x4+x6+x8 = 5540. Com esse valor, t = q+1/q é a raiz irracional de uma cúbica (só sai por Cardano), e a resposta é S $\approx$ 11196,5. Com 5440 sai tudo exato (q = 2, S = 10912). A nota depende de como a banca tratar isso: anular a questão (como sugere o Poliedro), aceitar 10912 ou aceitar a resposta literal. A resolução cobre os dois casos.%
\item%
\textbf{Q7} \textit{(key)} A resolução do gabarito afirma que 'o primo 1013 não divide 1025!$\cdot$1001!'. Isso é falso: 1013 $\leq$ 1025, então 1013 | 1025!. O argumento correto conta a multiplicidade: v\_1013(2026!) = 2 e v\_1013(1025!$\cdot$1001!) = 1, logo v\_1013 de C(2026,1001) é 1. Pode ser só um resumo malfeito, mas um corretor que siga o gabarito ao pé da letra aceitaria um argumento errado. Também poderia estranhar uma solução correta que use multiplicidade ou o lema de Gauss.%
\end{itemize}%
\section*{Questão 1}%
\href{https://poliedroresolve.sistemapoliedro.com.br/vestibulares/ime/2027/ime-2-fase-dia-1-26-10-2026-2-fase-dia-1/dissertativa-1-matematica-2-fase-dia-1-ime-2027}{Resolução do gabarito}\par\smallskip%
\begin{longtable}{@{}p{0.06\textwidth} p{0.33\textwidth} p{0.33\textwidth} p{0.1\textwidth} p{0.08\textwidth}@{}}\toprule Item & Resolução & Gabarito & Veredito & Certo \\ \midrule
Área de P ? I & 12$\surd$3 - 18 = 6(2$\surd$3 - 3) $\approx$ 2,78 u.a. & 12$\surd$3 - 18 ($\approx$ 2,785) & igual & ambos \\
\bottomrule\end{longtable}%
\begin{itemize}%
\item%
\textbf{Área de P ? I)} Conferi a conta e está certa. Cada hexágono tem área 3$\surd$3/2. A interseção é um dodecágono regular de apótema $\surd$3/2, com área 12$\cdot$(3/4)$\cdot$(2-$\surd$3) = 18 - 9$\surd$3 $\approx$ 2,412. Por inclusão{-}exclusão, 3$\surd$3 - (18 - 9$\surd$3) = 12$\surd$3 - 18 $\approx$ 2,7846. A conferência da resolução também bate: hexágono P mais 6 pontas de altura h = 1 - $\surd$3/2 e ângulo de 120$^\circ$ no vértice, cada uma com área $\surd$3h$^2$ = 7$\surd$3/4 - 3. Também confirmei com uma integração numérica em grade 2000$\times$2000, que deu $\approx$ 2,7846. A transcrição do enunciado está igual ao original.%
\end{itemize}%
\section*{Questão 2}%
\href{https://poliedroresolve.sistemapoliedro.com.br/vestibulares/ime/2027/ime-2-fase-dia-1-26-10-2026-2-fase-dia-1/dissertativa-2-matematica-2-fase-dia-1-ime-2027}{Resolução do gabarito}\par\smallskip%
\begin{longtable}{@{}p{0.06\textwidth} p{0.33\textwidth} p{0.33\textwidth} p{0.1\textwidth} p{0.08\textwidth}@{}}\toprule Item & Resolução & Gabarito & Veredito & Certo \\ \midrule
Questão 2 (probabilidade) & P = 2/5 (40\%). As raízes favoráveis são z = 1/2 $\pm$ ($\surd$3/2)i, que dão igualdade. & P = 2/5 & igual & ambos \\
\bottomrule\end{longtable}%
\begin{itemize}%
\item%
\textbf{Questão 2 (probabilidade))} Conferi o enunciado no PDF e a transcrição está fiel. As raízes são as raízes sextas da unidade diferentes de 1, isto é, cis(k$\pi$/3) com k=1..5, todas simples. Pela conta exata, |z-3/2|$^2$ = 13/4 - 3cos$\theta$, e a inequação equivale a cos$\theta$ $\geq$ 1/2. Isso vale só para k=1 e k=5, que ficam exatamente na fronteira (|z-3/2|$^2$ = 7/4). Para k=2 e k=4 o valor é 19/4, e para k=3 é 25/4. Conferi também numericamente em python. Dá 2 favoráveis em 5, então P = 2/5. A resolução segue o mesmo caminho do gabarito, com mais detalhe.%
\end{itemize}%
\paragraph{Pontos de dúvida}%
\begin{itemize}%
\item%
\textit{(statement)} A resolução e o gabarito supõem que as 5 raízes distintas são equiprováveis. Essa é a leitura natural de 'escolher ao acaso uma raiz qualquer' e não há multiplicidade que a complique, já que todas as raízes são simples.%
\item%
\textit{(solution)} As duas raízes favoráveis estão exatamente sobre a circunferência |z-3/2| = $\surd$7/2. O resultado depende de a desigualdade ser não estrita ($\leq$), como está no original. A resolução observa isso corretamente.%
\end{itemize}%
\section*{Questão 3}%
\href{https://poliedroresolve.sistemapoliedro.com.br/vestibulares/ime/2027/ime-2-fase-dia-1-26-10-2026-2-fase-dia-1/dissertativa-3-matematica-2-fase-dia-1-ime-2027}{Resolução do gabarito}\par\smallskip%
\begin{longtable}{@{}p{0.06\textwidth} p{0.33\textwidth} p{0.33\textwidth} p{0.1\textwidth} p{0.08\textwidth}@{}}\toprule Item & Resolução & Gabarito & Veredito & Certo \\ \midrule
Q3 --- menor inteiro x com f(x) + f(-1/x) $\geq$ 0,999 & x = 7 & x = 7 & igual & ambos \\
\bottomrule\end{longtable}%
\begin{itemize}%
\item%
\textbf{Q3 --- menor inteiro x com f(x) + f(-1/x) $\geq$ 0,999)} Conferi com o enunciado original (p. 1), que bate com a transcrição. A conta está certa. Para x $\geq$ 0 todos os termos entram, e f(x) = $\Sigma$(1/3)\^{}n = 1/2. Para x = -1/k, g\_n vale 1 só quando n < k, então f(-1/k) = (1/2)(1 - 3\^{}\{1-k\}). A soma dá 1 - 1/(2$\cdot$3\^{}\{k-1\}) $\geq$ 0,999, o que equivale a 3\^{}\{k-1\} $\geq$ 500. Como 3\^{}5 = 243 < 500 $\leq$ 729 = 3\^{}6, isso dá k $\geq$ 7. Para todo inteiro x $\leq$ -1, f(x) = 0 e f(-1/x) = 1/2, logo a soma é 1/2. Fiz a verificação exata em Python, com frações, para x de -5 a 9: x = 6 dá 0,99794 (não serve) e x = 7 dá 0,99931 (serve). A resolução descarta os inteiros negativos e x = 0 de forma explícita, e o gabarito apenas supõe x positivo sem justificar.%
\end{itemize}%
\paragraph{Pontos de dúvida}%
\begin{itemize}%
\item%
\textit{(key)} O gabarito parte de x $\in$ $\mathbb{Z}$+ sem justificar. O mais rigoroso é mostrar que todo inteiro negativo dá soma 1/2 < 0,999 e que x = 0 fica fora porque -1/x não está definido. Isso é necessário para afirmar que 7 é o menor inteiro, e não apenas o menor inteiro positivo. A resolução faz essa verificação, e um corretor rigoroso poderia cobrar esse passo.%
\item%
\textit{(statement)} O enunciado não diz explicitamente que x $\neq$ 0, mas f(-1/x) exige isso. As duas resoluções adotam essa hipótese, que é a única leitura razoável.%
\end{itemize}%
\section*{Questão 4}%
\href{https://poliedroresolve.sistemapoliedro.com.br/vestibulares/ime/2027/ime-2-fase-dia-1-26-10-2026-2-fase-dia-1/dissertativa-4-matematica-2-fase-dia-1-ime-2027}{Resolução do gabarito}\par\smallskip%
\begin{longtable}{@{}p{0.06\textwidth} p{0.33\textwidth} p{0.33\textwidth} p{0.1\textwidth} p{0.08\textwidth}@{}}\toprule Item & Resolução & Gabarito & Veredito & Certo \\ \midrule
Questão 4 & 1345 & 1345 & igual & ambos \\
\bottomrule\end{longtable}%
\begin{itemize}%
\item%
\textbf{Questão 4)} Refiz a conta e o resultado é 1345. Agrupando os fatores por x\_j, cada produto vale ?\_i(y\_i+x\_j) = -P(-x\_j) = x\_j$^3$ - 4x\_j + 2, com P(y) = y$^3$ - 4y - 2. Como x$^2$ = 7x + 3, temos x$^3$ = 7x$^2$ + 3x = 52x + 21, e o fator fica 48x\_j + 23. Pelas relações de Girard (x$_{1}$+x$_{2}$ = 7, x$_{1}$x$_{2}$ = -3): 2304$\cdot$(-3) + 1104$\cdot$7 + 529 = -6912 + 7728 + 529 = 1345. Calculei também em Python (numpy) o produto direto dos seis fatores com as raízes numéricas e deu 1345,000000\ldots{}, confirmando. O gabarito faz o mesmo caminho, só com o sinal posto em cada fator: (-48x$_{1}$-23)(-48x$_{2}$-23), que é o mesmo número. O enunciado transcrito confere com o PDF original.%
\end{itemize}%
\section*{Questão 5}%
\href{https://poliedroresolve.sistemapoliedro.com.br/vestibulares/ime/2027/ime-2-fase-dia-1-26-10-2026-2-fase-dia-1/dissertativa-5-matematica-2-fase-dia-1-ime-2027}{Resolução do gabarito}\par\smallskip%
\begin{longtable}{@{}p{0.06\textwidth} p{0.33\textwidth} p{0.33\textwidth} p{0.1\textwidth} p{0.08\textwidth}@{}}\toprule Item & Resolução & Gabarito & Veredito & Certo \\ \midrule
Questão 5 (x1+x3+x5+x7+x9) & Com o dado 5540: x5 = 512 e S = 512 + 5540t - 512t$^2$ $\approx$ 11196,5, onde t = q+1/q $\approx$ 2,5116 é a única raiz real de 128t$^3$ - 256t - 1385 = 0 (irracional). A resolução aponta provável erro de digitação: com 5440, q = 2, x\_k = 2\^{}(k+4) e S = 10912. & Com 5540 não há solução em forma fechada simples. Supondo erro de digitação (5440): q = 2, S = 10912. O gabarito sugere anular a questão. & igual & ambos \\
\bottomrule\end{longtable}%
\begin{itemize}%
\item%
\textbf{Questão 5 (x1+x3+x5+x7+x9))} Conferi o enunciado no PDF original: ele traz mesmo 5540, então o erro é da banca e não da transcrição. As duas respostas chegam a S = 10912 na hipótese de 5440 e concordam que, com 5540, a equação em t não tem raiz racional. Refiz as contas. A soma dos pares é 512(t$^3$-2t) e a dos ímpares é 512(t$^{4}$-3t$^2$+1), e as duas identidades estão certas. A raiz real de 128t$^3$-256t-1385 = 0 é t $\approx$ 2,5116 (u = 8t $\approx$ 20,0928), e a fórmula de Cardano dá o mesmo valor. Com isso, q $\approx$ 2,01543, r $\approx$ 1,0111 > 0 e S $\approx$ 11196,4994. A soma dos pares volta a dar 5540 e o produto dos ímpares dá 2\^{}45. A irracionalidade também se confirma: h(u) = u$^3$-128u-5540 vale -100 em u = 20 e 1033 em u = 21, e nenhum divisor de 5540 é raiz. Com 5440, 8t$^3$-16t-85 = (2t-5)(4t$^2$+10t+17), logo t = 5/2, q = 2 e S = 32+128+512+2048+8192 = 10912. A resolução ainda dá a resposta para o dado literal, o que o gabarito não faz. Isso completa o gabarito e não o contradiz.%
\end{itemize}%
\paragraph{Pontos de dúvida}%
\begin{itemize}%
\item%
\textit{(statement)} O enunciado oficial (PDF, pág. 2) traz x2+x4+x6+x8 = 5540. Com esse valor, t = q+1/q é a raiz irracional de uma cúbica (só sai por Cardano), e a resposta é S $\approx$ 11196,5. Com 5440 sai tudo exato (q = 2, S = 10912). A nota depende de como a banca tratar isso: anular a questão (como sugere o Poliedro), aceitar 10912 ou aceitar a resposta literal. A resolução cobre os dois casos.%
\item%
\textit{(key)} O gabarito diz que, com 5540, a questão 'não tem solução em forma fechada'. Isso não é exato: a solução existe, é única e tem forma fechada em radicais pela fórmula de Cardano, u = ?(2770 + (2/9)$\surd$153803361) + ?(2770 - (2/9)$\surd$153803361), com t = u/8. O que falta é uma forma fechada simples ou racional, como diz o próprio texto da resolução do gabarito.%
\item%
\textit{(solution)} A hipótese de que o valor correto seria 5440 é uma conjectura plausível (só um dígito trocado, e dá q = 2 inteiro). Outras correções do enunciado também seriam possíveis. A resolução marca isso corretamente como 'provável erro de digitação' e não troca a resposta literal.%
\end{itemize}%
\section*{Questão 6}%
\href{https://poliedroresolve.sistemapoliedro.com.br/vestibulares/ime/2027/ime-2-fase-dia-1-26-10-2026-2-fase-dia-1/dissertativa-6-matematica-2-fase-dia-1-ime-2027}{Resolução do gabarito}\par\smallskip%
\begin{longtable}{@{}p{0.06\textwidth} p{0.33\textwidth} p{0.33\textwidth} p{0.1\textwidth} p{0.08\textwidth}@{}}\toprule Item & Resolução & Gabarito & Veredito & Certo \\ \midrule
Questão 6 --- probabilidade de f\^{}(9)(1) = 1 & P = 1/3 (120960 permutações favoráveis de 362880) & P = 1/3 & igual & ambos \\
\bottomrule\end{longtable}%
\begin{itemize}%
\item%
\textbf{Questão 6 --- probabilidade de f\^{}(9)(1) = 1)} A resolução e o gabarito chegam a 1/3. Pelo enunciado original (p. 2), f\^{}(9)(1) = 1 vale exatamente quando o comprimento k do ciclo que contém o 1 divide 9, ou seja, k = 1, 3 ou 9. Para cada k há 8!/(9-k)! $\cdot$ (9-k)! = 8! permutações, então P = 3$\cdot$8!/9! = 1/3. A contagem do gabarito para o ciclo de comprimento 3 também está certa: C(8,2)$\cdot$2$\cdot$6! = 28$\cdot$2$\cdot$720 = 40320 = 8!. Testei por força bruta em Python as 9! = 362880 permutações. Exatamente 120960 delas satisfazem f\^{}(9)(1) = 1, o que dá 1/3.%
\end{itemize}%
\section*{Questão 7}%
\href{https://poliedroresolve.sistemapoliedro.com.br/vestibulares/ime/2027/ime-2-fase-dia-1-26-10-2026-2-fase-dia-1/dissertativa-7-matematica-2-fase-dia-1-ime-2027}{Resolução do gabarito}\par\smallskip%
\begin{longtable}{@{}p{0.06\textwidth} p{0.33\textwidth} p{0.33\textwidth} p{0.1\textwidth} p{0.08\textwidth}@{}}\toprule Item & Resolução & Gabarito & Veredito & Certo \\ \midrule
7 (demonstração) & tr(AB) = C(m+n+2, n+1). Para m = 1024 e n = 1000 isso dá C(2026, 1001). Como 1001$\cdot$C(2026,1001) = 2$\cdot$1013$\cdot$C(2025,1000) e mdc(1001,1013) = 1, segue que 1013 | tr(AB). A resolução confere ainda pela fórmula de Legendre: o expoente de 1013 é 1. & tr(AB) = C(m+n+2, n+1) = C(2026, 1001) = 2026!/(1025!$\cdot$1001!), divisível por 1013. & igual & ambos \\
\bottomrule\end{longtable}%
\begin{itemize}%
\item%
\textbf{7 (demonstração))} O enunciado confere com o original (pág. 2). A fórmula do traço é a mesma nas duas, obtida por Vandermonde em c11, c22 e c33 e depois Stifel duas vezes. Conferi em Python: tr(AB) = C(m+n+2, n+1) para todo 1 $\leq$ n < m $\leq$ 24, e também direto da definição para m = 1024 e n = 1000, com resultado igual a C(2026,1001) $\equiv$ 0 (mod 1013) e expoente de 1013 igual a 1. O caso de teste m = 2, n = 1 dá 10, como a resolução diz. O argumento da resolução está certo: 1001$\cdot$C(2026,1001) = 2026$\cdot$C(2025,1000) = 1013$\cdot$2$\cdot$C(2025,1000), e o lema de Gauss se aplica porque mdc(1001,1013) = mdc(1001,12) = 1. Já a justificativa escrita no gabarito diz que '1013 não divide 1025!$\cdot$1001!', o que é falso, já que 1013 $\leq$ 1025 e portanto 1013 | 1025!. O certo é: 1013 aparece duas vezes em 2026! (pelos fatores 1013 e 2026) e só uma vez em 1025!$\cdot$1001!, então sobra um fator 1013. A conclusão do gabarito está certa, mas esse passo, do jeito que está escrito, não.%
\end{itemize}%
\paragraph{Pontos de dúvida}%
\begin{itemize}%
\item%
\textit{(key)} A resolução do gabarito afirma que 'o primo 1013 não divide 1025!$\cdot$1001!'. Isso é falso: 1013 $\leq$ 1025, então 1013 | 1025!. O argumento correto conta a multiplicidade: v\_1013(2026!) = 2 e v\_1013(1025!$\cdot$1001!) = 1, logo v\_1013 de C(2026,1001) é 1. Pode ser só um resumo malfeito, mas um corretor que siga o gabarito ao pé da letra aceitaria um argumento errado. Também poderia estranhar uma solução correta que use multiplicidade ou o lema de Gauss.%
\item%
\textit{(solution)} No item T3, a observação 'n+1 $\leq$ m porque m>n' é desnecessária. Com a convenção C(a,c) = 0 para c > a, a identidade de Vandermonde vale de qualquer jeito. A observação não está errada e não prejudica a solução.%
\end{itemize}%
\section*{Questão 8}%
\href{https://poliedroresolve.sistemapoliedro.com.br/vestibulares/ime/2027/ime-2-fase-dia-1-26-10-2026-2-fase-dia-1/dissertativa-8-matematica-2-fase-dia-1-ime-2027}{Resolução do gabarito}\par\smallskip%
\begin{longtable}{@{}p{0.06\textwidth} p{0.33\textwidth} p{0.33\textwidth} p{0.1\textwidth} p{0.08\textwidth}@{}}\toprule Item & Resolução & Gabarito & Veredito & Certo \\ \midrule
Q8 (demonstração da identidade S\_n) & Pela prostaférese, cos $\alpha$ + cos(2k+1)$\alpha$ = 2cos(k$\alpha$)cos((k+1)$\alpha$). Como tg((k+1)$\alpha$) - tg(k$\alpha$) = sen $\alpha$/(cos(k$\alpha$)cos((k+1)$\alpha$)), a soma telescopa: S\_n = {[}tg((n+1)$\alpha$) - tg $\alpha${]}/(2 sen $\alpha$) = sen(n$\alpha$)/(sen(2$\alpha$)cos((n+1)$\alpha$)). & Demonstração: prostaférese, depois tg x - tg y = sen(x-y)/(cos x cos y) e soma telescópica: sen $\alpha$$\cdot$S\_n = (1/2)(tg((n+1)$\alpha$) - tg $\alpha$), logo S\_n = sen(n$\alpha$)/(sen 2$\alpha$$\cdot$cos((n+1)$\alpha$)). & igual & ambos \\
\bottomrule\end{longtable}%
\begin{itemize}%
\item%
\textbf{Q8 (demonstração da identidade S\_n))} O enunciado transcrito confere com a página 2 da prova. A resolução segue o mesmo caminho do gabarito: prostaférese, diferença de tangentes e soma telescópica. Refiz cada passo. A prostaférese está certa: (p+q)/2 = (k+1)$\alpha$ e (p-q)/2 = k$\alpha$. O numerador da diferença de tangentes é sen((k+1)$\alpha$ - k$\alpha$) = sen $\alpha$. Na simplificação final, sen((n+1)$\alpha$ - $\alpha$) = sen(n$\alpha$), e sen 2$\alpha$ = 2 sen $\alpha$ cos $\alpha$. A resolução também diz onde entra cada hipótese: cos(j$\alpha$) $\neq$ 0 para j = 1, \ldots{}, n+1 garante que os denominadores e as tangentes existem, e sen $\alpha$ $\neq$ 0 permite dividir por sen $\alpha$. Ela é até mais completa que o gabarito nesse ponto. O caso n = 1 confere. Uma checagem numérica com cerca de 20 mil pares (n, $\alpha$) aleatórios deu erro relativo máximo de cerca de 4$\cdot$10$^{-}$$^{1}$$^{0}$, tanto para a forma final quanto para a forma com tangentes.%
\end{itemize}%
\paragraph{Pontos de dúvida}%
\begin{itemize}%
\item%
\textit{(solution)} No fim do passo 4, a frase "Esse valor é não nulo, pois sen $\alpha$ $\neq$ 0 e cos $\alpha$ $\neq$ 0" pretende falar de sen(2$\alpha$), que está no denominador. Mas a redação também pode ser lida como se S\_n fosse não nulo, o que é falso em geral. Exemplo: com n = 3 e $\alpha$ = $\pi$/3 as hipóteses valem, e as parcelas -2, 1 e 1 somam S\_3 = 0. Isso não afeta a demonstração; basta escrever "sen(2$\alpha$) $\neq$ 0" de forma explícita.%
\item%
\textit{(solution)} Escorregão de nome: no passo 4 a resolução diz "subtração de senos", mas a fórmula usada é a do seno da diferença, sen(a-b). Não afeta o conteúdo.%
\end{itemize}%
\section*{Questão 9}%
\href{https://poliedroresolve.sistemapoliedro.com.br/vestibulares/ime/2027/ime-2-fase-dia-1-26-10-2026-2-fase-dia-1/dissertativa-9-matematica-2-fase-dia-1-ime-2027}{Resolução do gabarito}\par\smallskip%
\begin{longtable}{@{}p{0.06\textwidth} p{0.33\textwidth} p{0.33\textwidth} p{0.1\textwidth} p{0.08\textwidth}@{}}\toprule Item & Resolução & Gabarito & Veredito & Certo \\ \midrule
Lugar geométrico de M & Elipse 4x$^2$/49 + y$^2$ = 1 (4x$^2$ + 49y$^2$ = 49), com semieixos 7/2 e 1, sem os vértices ($\pm$7/2, 0) & 4x$^2$/49 + y$^2$ = 1 (elipse) & igual & ambos \\
Pontos de interseção com x = $\pm$c & c = 2$\surd$3; pontos (2$\surd$3, $\pm$1/7) e (-2$\surd$3, $\pm$1/7) & (2$\surd$3, 1/7), (2$\surd$3, -1/7), (-2$\surd$3, 1/7), (-2$\surd$3, -1/7) & igual & ambos \\
\bottomrule\end{longtable}%
\begin{itemize}%
\item%
\textbf{Lugar geométrico de M)} Refiz a conta: a normal em P=(x0,y0) tem direção (x0, 4y0) e corta y=0 em Q=(3x0/4, 0). Então M=(7x0/8, y0/2). Substituindo x0=8X/7 e y0=2Y em x0$^2$+4y0$^2$=16, sai 4X$^2$/49+Y$^2$=1. Conferi numericamente para vários $\theta$ com P=(4cos$\theta$, 2sin$\theta$). A resolução ainda exclui os vértices ($\pm$7/2, 0), que viriam de P=($\pm$4,0). Isso é um refinamento que o gabarito não traz e não muda o resultado.%
\item%
\textbf{Pontos de interseção com x = $\pm$c)} c = $\surd$(16-4) = 2$\surd$3. Com x$^2$ = 12 no lugar: y$^2$ = 1 - 48/49 = 1/49, logo y = $\pm$1/7. Como 2$\surd$3 $\approx$ 3,46 < 7/2, cada reta dá dois pontos. Testei a recíproca com M=(2$\surd$3, 1/7): P=(16$\surd$3/7, 2/7) está na elipse (768/49 + 16/49 = 16), Q=(12$\surd$3/7, 0) e o ponto médio é (2$\surd$3, 1/7). Os quatro pontos coincidem com o gabarito.%
\end{itemize}%
\paragraph{Pontos de dúvida}%
\begin{itemize}%
\item%
\textit{(solution)} Caso degenerado P=($\pm$4,0): a normal é o próprio eixo x e Q não fica determinado. A resolução exclui esses pontos e, com eles, os vértices ($\pm$7/2, 0) do lugar. É uma convenção razoável e o gabarito declara que não trata o caso. Como c = 2$\surd$3 $\neq$ 7/2, isso não afeta os pontos pedidos.%
\item%
\textit{(solution)} Na verificação aparece a forma não simplificada (28$\surd$3/14, 1/7). Ela é igual a (2$\surd$3, 1/7) e não há erro, é só a apresentação.%
\end{itemize}%
\section*{Questão 10}%
\href{https://poliedroresolve.sistemapoliedro.com.br/vestibulares/ime/2027/ime-2-fase-dia-1-26-10-2026-2-fase-dia-1/dissertativa-10-matematica-2-fase-dia-1-ime-2027}{Resolução do gabarito}\par\smallskip%
\begin{longtable}{@{}p{0.06\textwidth} p{0.33\textwidth} p{0.33\textwidth} p{0.1\textwidth} p{0.08\textwidth}@{}}\toprule Item & Resolução & Gabarito & Veredito & Certo \\ \midrule
a & V\_max = 192$\surd$3 cm$^3$ $\approx$ 332,6 cm$^3$ & V\_max = 192$\surd$3 cm$^3$ ($\approx$ 332,55 cm$^3$) & igual & ambos \\
b & Altura 4$\surd$3 cm; base quadrada de lado 4$\surd$3 cm $\times$ 4$\surd$3 cm (cubo de aresta 4$\surd$3 cm inscrito na esfera) & Cubo: altura e lados da base iguais a 4$\surd$3 cm & igual & ambos \\
\bottomrule\end{longtable}%
\begin{itemize}%
\item%
\textbf{a)} As duas dão o mesmo valor. Refiz a conta: f(m) = 4m(36 - m$^2$) tem máximo em m = 2$\surd$3, e f = 4$\cdot$2$\surd$3$\cdot$24 = 192$\surd$3 $\approx$ 332,554. Pelo MA--MG do gabarito, 48\^{}(3/2) = 192$\surd$3. Uma busca numérica aleatória também ficou em $\approx$ 332,55. A resolução é mais geral que o gabarito: vale para qualquer base quadrilátera simples, convexa ou não, usando S = $\frac12$|AC$\times$BD| $\leq$ $\frac12$d$_{1}$d$_{2}$ $\leq$ 2$\rho$$^2$. Também não supõe que os 8 vértices estejam na esfera: usa só que o prisma está dentro da bola, com $\rho$$^2$ = 36 - m$^2$ e h $\leq$ 2m. O gabarito supõe uma base retangular inscrita e h = $\surd$(144 - a$^2$ - b$^2$) sem justificar. O resultado não muda.%
\item%
\textbf{b)} As duas respostas coincidem. Conferi: com m = 2$\surd$3, h = 2m = 4$\surd$3 e $\rho$$^2$ = 24, o lado do quadrado é $\rho$$\surd$2 = 4$\surd$3. Cada vértice fica a uma distância$^2$ de 24 + 12 = 36 do centro, e a diagonal do cubo vale 4$\surd$3$\cdot$$\surd$3 = 12 = 2R. A resolução também verifica que a condição do enunciado (os quatro vértices de uma face na superfície) é satisfeita pelo cubo.%
\end{itemize}%
\paragraph{Pontos de dúvida}%
\begin{itemize}%
\item%
\textit{(key)} A resolução do gabarito supõe, sem justificar, que a base é retangular e que os 8 vértices estão na esfera (h = $\surd$(144 - a$^2$ - b$^2$)). O enunciado só pede um 'prisma quadrangular reto' contido na esfera, com uma face tendo os vértices na superfície. A resposta final está certa, mas a justificativa é incompleta. A resolução conferida cobre o caso geral: qualquer quadrilátero simples como base e prisma apenas contido na bola.%
\item%
\textit{(key)} Erro de digitação na fonte do gabarito: aparece 'a = 4$\surd$(3 cm)' no lugar de a = 4$\surd$3 cm. Não muda a resposta.%
\item%
\textit{(statement)} A condição 'os quatro vértices de uma das faces pertencem à superfície' não interfere no máximo: a cota V $\leq$ 192$\surd$3 vale para qualquer prisma reto contido na bola, e o cubo ótimo tem todos os 8 vértices na esfera. Se a intenção fosse 'exatamente uma face', o supremo não seria atingido e não haveria máximo. A leitura natural ('pelo menos uma face'), adotada pela resolução e pelo gabarito, é a coerente.%
\end{itemize}%
\end{document}